\solution{Қара пакет}{15.0}

Барлық өлшеулерді батареяның көмегімен жүргіземіз, яғни қара жәшікті сол арқылы зарядтаймыз. Сондықтан оның ЭҚК-ін өлшейміз: $V_0 = \qty{8.89}{\V}$.

Одан әрі қара жәшікті батарея мен амперметрге тізбектей жалғаймыз. Қара жәшік арқылы өтетін ток күшін өлшейміз: $I_r = \qty{1.81}{\mA}$. Батареяның полярлығын ауыстырып, бұл жағдайда да ток күшін өлшеп көрейік. Ток тұрақталған сәтті күте отырып, ток күшінің мынадай мәнге ие болғанын көреміз: $I_l = \qty{0.06}{\mA}$. Ток күшінің әртүрлі бағыттарда әртүрлі болу фактісі тізбекте диодтың бар екенін көрсетеді. Бұл ток мәні нөлге өте жақын және салыстырмалы түрде айтарлықтай аз. Сондықтан мұндай полярлық кезінде тізбекте ток жүрмейді және диод резистормен тізбектей жалғанған деп есептей аламыз. Осылайша, қара жәшіктің бір тармағы — тізбектей жалғанған диод пен резистор. Бұл ретте шығыстарға жалғанған басқа резисторлар жоқ, өйткені екінші жағдайда ток нөлге тең болды.

Ток мәнінен бірінші резистордың кедергісін оңай таба аламыз:
\begin{eq}
    R_1 = \frac{V_0}{I_r} \approx \qty{4.91}{\kOhm}
\end{eq}

Амперметрді жалғау кезінде токтың уақыт өте келе өзгергенін және біраз уақыттан кейін ғана стационарлық мәнге жеткенін байқауға болады. Мұндай құбылыс конденсатор болған жағдайда орын алады. Бұдан бөлек, екі бағыттағы ток көрсеткіштерінен қара жәшіктің шығыстарына диоды бар резисторлар немесе басқа элементтермен тізбектей жалғанған конденсаторлар қосылған деген қорытынды жасауға болады. Бір типті бірнеше параллель жалғанған тармақтар өзін ұқсас ұстайтындықтан, олар тек екі тармағы бар ең қарапайым нұсқадағы схемаға эквивалентті болады: бірі — резистор мен диод, екіншісі — басқа элементтері бар конденсатор.

Конденсаторға байланысты тізбектің өзін-өзі ұстауын зерттеу үшін келесідей істейміз: Қара жәшік, батарея және вольтметрді параллель жалғаймыз. Біраз уақыт күтеміз (амперметрді жалғаған кезде ток тұрақталатындай уақыт шамасы) және батареяны тізбектен ажыратамыз. Одан кейін кернеудің уақытқа тәуелділігін зерттейтін боламыз.

Бізге белгілі болғандай, қара жәшіктің схемасын сызайық. Сондай-ақ, нені өлшеп жатқанымызды түсіну үшін вольтметрді де көрсетейік:

\cpic{0.4}{figures/test.jpg}

Схемадан біз өлшеп жатқан кернеудің резистордағы кернеудің түсуі екені айқынырақ бола түседі.

\begin{table}[h!]
\centering
\caption{Разрядталу деректері (1-бөлім)}
\resizebox{\textwidth}{!}{
\begin{tabular}{l|*{16}{c}}
\hline
$V$, В & 1.263 & 1.096 & 0.923 & 0.861 & 0.548 & 0.490 & 0.448 & 0.390 &0.371& 0.354 & 0.338 & 0.326 & 0.318 & 0.304 & 0.296 & 0.286 \\
$t$, с & 2.07 & 2.39 & 2.94 & 3.39 & 5.23 & 5.96 & 6.97 & 8.45 &9.49& 10.47 & 11.47 & 12.81 & 14.09 & 15.61 & 17.26 & 18.37 \\
$\ln V$ & 0.233 & 0.092 & -0.080 & -0.150 & -0.601 & -0.713 & -0.803 & -0.942 &-0.992&  -1.038 & -1.085 & -1.121 & -1.146 & -1.191 & -1.217 & -1.252 \\
\hline
\end{tabular}
}
\end{table}

Осы тәуелділіктің графигін сызайық:

\begin{table}[h!]
\centering
\caption{Разрядталу деректері (2-бөлім)}
\resizebox{\textwidth}{!}{
\begin{tabular}{l|*{14}{c}}
\hline
$V$, В & 0.279 & 0.265 & 0.258 & 0.246 & 0.234 & 0.227 & 0.220 & 0.215 & 0.210 & 0.206 & 0.197 & 0.192 & 0.187 & 0.179 \\
$t$, с & 20.47 & 23.03 & 26.58 & 30.23 & 36.41 & 41.07 & 45.62 & 50.65 & 54.99 & 59.57 & 65.20 & 75.23 & 82.92 & 108.51 \\
$\ln V$ & -1.277 & -1.328 & -1.355 & -1.402 & -1.452 & -1.483 & -1.514 & -1.537 & -1.561 & -1.580 & -1.625 & -1.650 & -1.677 & -1.720 \\
\hline
\end{tabular}
}
\end{table}

\vspace{1cm}

\section*{Графиктер}

\begin{figure}[h!]
    \centering
    % V(t) графигі
    \begin{tikzpicture}
    \begin{axis}[
        title={Кернеудің уақытқа тәуелділігі $V(t)$},
        xlabel={Уақыт $t$, с},
        ylabel={Кернеу $V$, В},
        grid=major,
        width=12cm,
        height=8cm,
        mark size=1.5pt
    ]
    \addplot[
        only marks,
        mark=*,
        color=blue
    ] coordinates {
        (2.07, 1.263) (2.39, 1.096) (2.94, 0.923) (3.39, 0.861) (5.23, 0.548) (5.96, 0.490) (6.97, 0.448) (8.45, 0.390) (10.47, 0.354) (11.47, 0.338) (12.81, 0.326) (14.09, 0.318) (15.61, 0.304) (17.26, 0.296) (18.37, 0.286) %(20.47, 0.279) (23.03, 0.265) (26.58, 0.258) (30.23, 0.246) (36.41, 0.234) (41.07, 0.227) (45.62, 0.220) (50.65, 0.215) (54.99, 0.210) (59.57, 0.206) (65.20, 0.197) (75.23, 0.192) (82.92, 0.187) (108.51, 0.179)
    };
    \end{axis}
    \end{tikzpicture}
\end{figure}


\begin{figure}[h!]
    \centering
    % V(t) графигі
    \begin{tikzpicture}
    \begin{axis}[
        title={Кернеудің уақытқа тәуелділігі $V(t)$},
        xlabel={Уақыт $t$, с},
        ylabel={Кернеу $V$, В},
        grid=major,
        width=12cm,
        height=8cm,
        mark size=1.5pt
    ]
    \addplot[
        only marks,
        mark=*,
        color=blue
    ] coordinates {
        (20.47, 0.279) (23.03, 0.265) (26.58, 0.258) (30.23, 0.246) (36.41, 0.234) (41.07, 0.227) (45.62, 0.220) (50.65, 0.215) (54.99, 0.210) (59.57, 0.206) (65.20, 0.197) (75.23, 0.192) (82.92, 0.187) (108.51, 0.179)
    };
    \end{axis}
    \end{tikzpicture}
\end{figure}

Кернеу графиктері экспоненциалдық тәуелділікке өте ұқсас, әсіресе процестің басталуын көрсететін бірінші график. Сондықтан осы бөлікті қарастырамыз. Бұл нүктелер үшін графикті $\ln V, t$ координаттарында сызамыз.

\begin{figure}[h!]
    \centering
    \begin{tikzpicture}
    \begin{axis}[
        title={$\ln V$-ның $t$-ға тәуелді жартылай логарифмдік графигі},
        xlabel={Уақыт $t$, с},
        ylabel={$\ln V$},
        grid=major,
        width=12cm,
        height=8cm,
        mark size=1.5pt
    ]

    % 2. Бастапқы бөлікке арналған тренд сызығы (2-6 с)
    \addplot[
        domain=1.5:6.5,
        samples=2,
        color=red,
        thick,
        dashed
    ] {-0.238*x + 0.668};
    \addlegendentry{Аппроксимация}

    \addplot[
        only marks,
        mark=triangle*,
        color=purple
    ] coordinates {
        (2.07, 0.233) (2.39, 0.092) (2.94, -0.080) (3.39, -0.150) (5.23, -0.601) (5.96, -0.713) (6.97, -0.803) (8.45, -0.942)(10.47, -1.038) (11.47, -1.085) (12.81, -1.121) %(14.09, -1.146) (15.61, -1.191) (17.26, -1.217) (18.37, -1.252) (20.47, -1.277) (23.03, -1.328) (26.58, -1.355) (30.23, -1.402) (36.41, -1.452) (41.07, -1.483) (45.62, -1.514) (50.65, -1.537) (54.99, -1.561) (59.57, -1.580) (65.20, -1.625) (75.23, -1.650) (82.92, -1.677) (108.51, -1.720)
    };
    \end{axis}
    \end{tikzpicture}
\end{figure}

График сызықтық болмай шықты. Дегенмен, әсіресе процестің басында сызықтық бөліктер бар. Одан кейінгі сызықтықтан ауытқу диодтың идеал еместігінен туындаған деп есептеледі.

Түзу аймақты $\ln V = -kt + b$ сызықтық функциясымен аппроксимациялаймыз. ЕККӘ (Ең кіші квадраттар әдісі) бойынша функция параметрлерін табамыз: $k \approx -0.238 \unit{\per\s}; b = 0.668 $.

Сипаттамалық уақыт $\tau$ мына өрнекке тең: $\tau = C\cdot (R_1 + R_x)$, мұндағы $R_x$ — белгісіз $X$ бөлігінің кедергісі. Шын мәнінде, қара жәшіктің өзін-өзі ұстауына тек $X$ элементінің кедергісі ғана әсер етеді. Барлық басқа схемалар эквивалентті болады, өйткені қалған элементтердің болуы құбылысқа әсер етпейді. Осылайша, соңғы схема келесідей болады:

\cpic{0.4}{figures/black_box.jpg}

Бұл жинақта шығыстағы кернеудің тәуелділігі мынадай түрге ие:
\begin{eq}
    V(t) = V_0 \frac{R_1}{R_1+R_2}\cdot \exp{(-t/\tau)} \implies \ln V = \ln V_0 \frac{R_1}{R_1+R_2} - \frac t \tau.
\end{eq}

Алдымен екінші резистордың кедергісін табайық:
\begin{eq}
    R_1 \cdot \left(\frac{V_0}{e^b} - 1\right) \approx\qty{17.4}{\kOhm}
\end{eq}

Енді кедергі мәндерін біле отырып, конденсатор сыйымдылығының мәнін таба аламыз: $C = \qty{188}{\uF}$.
